\chapter{逐点收敛与一致收敛}
\label{chap:pointwise-uniform-convergence}

数列的每一项是一个数，函数列的每一项却是一整族函数值。若只把自变量
固定下来，函数列便退化为普通数列；这种逐点考察自然给出逐点收敛。然而，
连续性、积分和微分都同时比较许多自变量处的函数值。要把极限过程送入这些
运算，通常需要一种对整个定义域同时有效的尾部控制。本章先区分这两种控制，
再把一致收敛写成有界函数空间中的范数收敛。

除非另有说明，本章的函数都取实值，公共定义域 \(E\) 是一个非空集合。

\section{函数列和逐点收敛}

函数列是映射序列
\[
 f_n:E\to\R\qquad(n\in\Npos).
\]
对固定的 \(x\in E\)，数列 \((f_n(x))\) 可以按照第一卷的方法讨论。

\begin{definition}[逐点收敛]
若对每个 \(x\in E\)，数列 \((f_n(x))\) 都收敛到 \(f(x)\)，则称
\((f_n)\) 在 \(E\) 上\term{逐点收敛}于 \(f\)，记作
\[
 f_n(x)\longrightarrow f(x)\quad(x\in E).
\]
展开量词，这表示
\[
 \forall x\in E\ \forall\varepsilon>0\ \exists N=N(x,\varepsilon)
 \ \forall n\ge N,
 \qquad |f_n(x)-f(x)|<\varepsilon.
\]
\end{definition}

这里的 \(N\) 可以随 \(x\) 改变。若某些点附近的收敛越来越慢，逐点定义
仍然允许为这些点另选更大的指标。

逐点极限若存在便是唯一的，因为每个固定 \(x\) 处的数列极限唯一。四则
运算也可以逐点进行：若 \(f_n\to f\)、\(g_n\to g\) 逐点成立，则
\(f_n+g_n\to f+g\)、\(f_ng_n\to fg\) 逐点成立；在 \(g(x)\ne0\) 的
点，\(g_n(x)\) 最终不为零，并有 \(f_n(x)/g_n(x)\to f(x)/g(x)\)。
仅有分母各项非零还不够，例如 \(f_n\equiv1\)、\(g_n\equiv1/n\)
的商没有有限极限。这些结论只是数列极限法则逐点使用，并没有给出不同点之间的
共同估计。

\begin{example}
在 \([0,1]\) 上令 \(f_n(x)=x^n\)。若 \(0\le x<1\)，则 \(x^n\to0\)；
而 \(f_n(1)=1\)。故逐点极限为
\[
 f(x)=
 \begin{cases}
 0,&0\le x<1,\\
 1,&x=1.
 \end{cases}
\]
每个 \(f_n\) 都连续，极限函数却在 \(1\) 处不连续。连续性不能仅由逐点
收敛传给极限。
\end{example}

逐点收敛还可以只在某个子集上考察。上例在每个 \([0,r]\)、\(0<r<1\)
上都有估计
\[
 \sup_{0\le x\le r}x^n=r^n\to0.
\]
端点 \(1\) 附近才是收敛失去共同速度的地方。

\section{一致收敛及量词顺序}

\begin{definition}[一致收敛]
若对任意 \(\varepsilon>0\)，存在只依赖于 \(\varepsilon\) 的
\(N\)，使对所有 \(x\in E\) 和所有 \(n\ge N\)，都有
\[
 |f_n(x)-f(x)|<\varepsilon,
\]
则称 \((f_n)\) 在 \(E\) 上\term{一致收敛}于 \(f\)，记作
\[
 f_n\rightrightarrows f\quad\text{于 }E.
\]
\end{definition}

与逐点收敛相比，变化只发生在两个量词的位置：
\[
\begin{array}{ll}
\text{逐点收敛：}&
 \forall x\ \forall\varepsilon\ \exists N\ \forall n\ge N,\\[0.3em]
\text{一致收敛：}&
 \forall\varepsilon\ \exists N\ \forall x\ \forall n\ge N.
\end{array}
\]
一致收敛要求同一个 \(N\) 控制定义域中所有点，因而必然推出逐点收敛。
反向推论由 \(x^n\) 的例子否定。

一致收敛的否定形式在构造反例时很有用：
\[
 f_n\not\rightrightarrows f
\]
当且仅当存在 \(\varepsilon_0>0\)，使对每个 \(N\)，都能找到
\(n\ge N\) 和 \(x_n\in E\)，满足
\[
 |f_n(x_n)-f(x_n)|\ge\varepsilon_0.
\]
这里的坏点允许随 \(n\) 移动。对 \([0,1]\) 上的 \(f_n(x)=x^n\)，取
\(x_n=2^{-1/n}<1\)，则 \(f(x_n)=0\)，而 \(f_n(x_n)=1/2\)，误差
始终为 \(1/2\)。端点 \(1\) 自身的误差反而恒为零；反例中的大误差
位于越来越靠近端点的左侧。删去端点后，同一组点仍然适用。

\begin{proposition}[一致收敛的限制与有限拼接]
若 \(f_n\rightrightarrows f\) 于 \(E\)，则在任意 \(A\subseteq E\)
上仍一致收敛。反之，若
\[
 E=A_1\cup\cdots\cup A_m
\]
是有限并，并且 \(f_n\rightrightarrows f\) 于每个 \(A_j\)，则
\(f_n\rightrightarrows f\) 于 \(E\)。
\end{proposition}

\begin{proof}
限制结论直接沿用原来的 \(N\)。对有限拼接，给定 \(\varepsilon>0\)，
分别取每个 \(A_j\) 所需的 \(N_j\)，令
\(N=\max_{1\le j\le m}N_j\)。任意 \(x\in E\) 至少属于一个
\(A_j\)，因此同一个 \(N\) 对整个 \(E\) 有效。
\end{proof}

“有限”不能直接删去。函数列可以在每个 \(A_m\) 上一致收敛，而所需指标
随 \(m\) 增大，以致在可数并上没有共同上界。例如 \(f_n(x)=x/n\) 在
每个有界区间上一致趋于零，在 \(\R\) 上却不一致收敛。

\section{上确界范数与一致 Cauchy 准则}

若误差函数 \(f_n-f\) 有界，定义
\[
 \norm{f_n-f}_\infty
   =\sup_{x\in E}|f_n(x)-f(x)|.
\]
下标 \(\infty\) 表示取所有点中最大的误差尺度，并非变量趋于无穷。

\begin{theorem}[上确界判据]
若 \(f_n-f\) 在 \(E\) 上有界，则
\[
 f_n\rightrightarrows f\quad\Longleftrightarrow\quad
 \norm{f_n-f}_\infty\to0.
\]
更一般地，允许上确界取值 \(+\infty\) 时，右端仍刻画一致收敛。
\end{theorem}

\begin{proof}
若上确界趋于零，则充分大时
\[
 |f_n(x)-f(x)|\le\norm{f_n-f}_\infty<\varepsilon
\]
对所有 \(x\) 同时成立。反之，若一致收敛，给定 \(\varepsilon>0\)，
充分大的 \(n\) 满足
\[
 |f_n(x)-f(x)|<\varepsilon/2\qquad(x\in E).
\]
于是 \(\norm{f_n-f}_\infty\le\varepsilon/2<\varepsilon\)。使用
\(\varepsilon/2\) 是为了避免从逐点的严格不等式直接推出上确界也严格小于
同一个 \(\varepsilon\)。
\end{proof}

这个判据把“一致”压缩成一个实数列的收敛。许多一致估计因而可以直接
使用三角不等式。例如
\[
 \norm{(f_n+g_n)-(f+g)}_\infty
 \le\norm{f_n-f}_\infty+\norm{g_n-g}_\infty.
\]

未知极限函数时，可以把数列的 Cauchy 思想用于函数列。

\begin{definition}[一致 Cauchy 条件]
称 \((f_n)\) 在 \(E\) 上满足一致 Cauchy 条件，如果
\[
 \forall\varepsilon>0\ \exists N\ \forall m,n\ge N\ \forall x\in E,
 \qquad |f_m(x)-f_n(x)|<\varepsilon.
\]
\end{definition}

\begin{theorem}[一致收敛的 Cauchy 准则]
实值函数列在 \(E\) 上一致收敛，当且仅当它满足一致 Cauchy 条件。
\label{thm:uniform-cauchy-criterion}
\end{theorem}

\begin{proof}
若 \(f_n\rightrightarrows f\)，取 \(N\) 使
\(|f_k(x)-f(x)|<\varepsilon/2\) 对所有 \(k\ge N\)、\(x\in E\)
成立，则
\[
 |f_m(x)-f_n(x)|<\varepsilon.
\]

反之，固定 \(x\in E\)。一致 Cauchy 条件说明 \((f_n(x))\) 是实
Cauchy 列，由实数完备性存在极限，记为 \(f(x)\)。给定
\(\varepsilon>0\)，取 \(N\) 使 \(m,n\ge N\) 时
\[
 |f_m(x)-f_n(x)|<\varepsilon/2
\]
对所有 \(x\) 成立。固定 \(n\ge N\) 和 \(x\)，令 \(m\to\infty\)，
由绝对值的连续性得到
\[
 |f(x)-f_n(x)|\le\varepsilon/2<\varepsilon.
\]
同一个 \(N\) 与 \(x\) 无关，故收敛是一致的。
\end{proof}

充分性证明使用了实数的 Cauchy 完备性。若值域换成不完备的度量空间，
一致 Cauchy 条件仍能使每个点处成为 Cauchy 列，却未必在原值域中得到
极限。这里已经出现了以后函数空间完备性的基本模式。

\section{有界函数空间的完备性}

记
\[
 \mathcal B(E)=\{f:E\to\R:f\text{ 有界}\}.
\]
逐点加法和数乘使它成为实线性空间。对 \(f\in\mathcal B(E)\)，定义
\[
 \norm f_\infty=\sup_{x\in E}|f(x)|.
\]

\begin{proposition}
上确界范数满足
\[
 \norm f_\infty\ge0,\qquad
 \norm f_\infty=0\Longleftrightarrow f=0,
\]
\[
 \norm{\lambda f}_\infty=|\lambda|\norm f_\infty,\qquad
 \norm{f+g}_\infty\le\norm f_\infty+\norm g_\infty.
\]
\end{proposition}

\begin{proof}
前三条来自绝对值的相应性质。对任意 \(x\in E\)，
\[
 |f(x)+g(x)|\le |f(x)|+|g(x)|
 \le\norm f_\infty+\norm g_\infty.
\]
对 \(x\) 取上确界便得三角不等式。
\end{proof}

因此 \(\norm{f-g}_\infty\) 给出有界函数之间的距离。上一节的上确界
判据说，一致收敛正是这个距离意义下的收敛。

\begin{theorem}
\label{thm:bounded-functions-complete}
\(\mathcal B(E)\) 在上确界范数下完备：若
\[
 \norm{f_m-f_n}_\infty\to0\qquad(m,n\to\infty),
\]
则存在 \(f\in\mathcal B(E)\)，使
\(\norm{f_n-f}_\infty\to0\)。
\end{theorem}

\begin{proof}
范数 Cauchy 条件与一致 Cauchy 条件等价。由
\theoremref{thm:uniform-cauchy-criterion}，存在函数 \(f:E\to\R\)
使 \(f_n\rightrightarrows f\)。还须验证 \(f\) 有界。取 \(N\) 使
\(\norm{f-f_N}_\infty<1\)，则对所有 \(x\in E\)，
\[
 |f(x)|\le |f(x)-f_N(x)|+|f_N(x)|
 \le1+\norm{f_N}_\infty.
\]
所以 \(f\in\mathcal B(E)\)，而上确界判据给出范数收敛。
\end{proof}

若 \(E\) 是区间，连续有界函数构成 \(\mathcal B(E)\) 的子空间。它是否
在上确界范数下仍然完备，取决于一致极限是否保持连续；下一章将证明答案
是肯定的。若把上确界范数换成积分平均意义的距离，收敛概念和完备化后的
函数空间都会改变，这一现象将在 Lebesgue 理论中系统出现。

\section{移动峰与尖峰函数}

逐点收敛失败于一致性的典型原因，是最大误差所在的点随 \(n\) 移动。
在 \([0,1]\) 上令
\[
 p_n(x)=\max\{1-3n|x-1/n|,0\}.
\]
其支集包含在
\([2/(3n),4/(3n)]\) 中，且 \(p_n(1/n)=1\)。固定 \(x>0\) 时，
这个支集最终落在 \(x\) 左侧；在 \(x=0\) 处函数值恒为零。因此
\[
 p_n(x)\to0\qquad(0\le x\le1),
\]
但
\[
 \norm{p_n}_\infty=1.
\]
图形的峰不断向原点移动，每一个固定点最终都躲开了峰，而整个区间始终
存在误差为 \(1\) 的点。

无界定义域还会产生另一种现象。对 \(x\ge0\) 令
\[
 q_n(x)=\frac{x}{n+x}.
\]
每个固定 \(x\) 都满足 \(q_n(x)\to0\)，但
\[
 \sup_{x\ge0}q_n(x)=1.
\]
这里坏点逃向无穷远处。限制在任意 \([0,M]\) 上则有
\[
 0\le q_n(x)\le \frac{M}{n+M}\to0,
\]
所以收敛在每个有界闭区间上一致，却不在整个半轴上一致。

这些例子说明，逐点收敛只观察固定位置的长期行为；一致收敛还要控制
误差峰的位置和高度。峰的宽度或面积是否趋零并不能单独决定一致收敛，
因为上确界范数只记录峰高。积分则会同时受到高度和宽度影响，这正是后续
交换极限与积分时需要分别讨论的地方。

\section*{习题}
\addcontentsline{toc}{section}{习题}

\subsection*{A组　定义与判别}
\begin{enumerate}[label=\textbf{\thechapter.\arabic*.}]
 \exerciseitem{pointwise-uniform-quantifier-negation}
 写出“\(f_n\) 不在 \(E\) 上一致收敛于 \(f\)”的完整量词形式，并说明
 为什么只找到一个固定点 \(x_0\) 使 \(f_n(x_0)\not\to f(x_0)\) 是更强的
 失败情形。
 \exerciseitem{uniform-power-on-truncated-interval}
 证明 \(x^n\rightrightarrows0\) 于 \([0,r]\)（\(0<r<1\)），并证明它在
 \([0,1)\) 上不一致收敛。
 \exerciseitem{uniform-rational-family}
 讨论 \(f_n(x)=x/(1+nx)\) 在 \([0,1]\) 上的逐点与一致收敛。
 \exerciseitem{uniform-reciprocal-tail}
 讨论 \(f_n(x)=1/(1+nx)\) 在 \([0,1]\) 和 \([a,1]\)（\(0<a\le1\)）上的
 逐点与一致收敛。
 \exerciseitem{uniform-finite-union}
 证明正文中的有限拼接命题，并构造一个可数并上失败的例子。
 \exerciseitem{supremum-criterion-strict-inequality}
 解释在上确界判据的反向证明中为什么使用 \(\varepsilon/2\) 可以避免
 严格不等号问题；举例说明“对每个 \(x\)，\(h(x)<1\)”不能推出
 \(\sup_Eh<1\)。
\end{enumerate}

\subsection*{B组　一致 Cauchy 条件与范数}
\begin{enumerate}[label=\textbf{\thechapter.\arabic*.},resume]
 \exerciseitem{uniform-cauchy-geometric-sums}
 不先写出极限函数，用一致 Cauchy 准则证明
 \(\sum_{k=0}^n x^k\) 在 \([0,r]\)、\(0<r<1\) 上一致收敛。
 \exerciseitem{uniform-cauchy-pointwise-insufficient}
 给出一个函数列，使每个 \(x\) 处的数列都是 Cauchy 列，但函数列不满足
 一致 Cauchy 条件。
 \exerciseitem{sup-norm-basic-identities}
 对 \(f,g\in\mathcal B(E)\) 证明
 \[
  |\norm f_\infty-\norm g_\infty|\le\norm{f-g}_\infty,
 \]
 并推出范数函数关于上确界距离是连续的。
 \exerciseitem{uniform-limit-boundedness}
 若每个 \(f_n\) 有界且 \(f_n\rightrightarrows f\)，证明 \(f\) 有界。
 再说明逐点极限不具有这一性质。
 \exerciseitem{bounded-functions-closed-subspace}
 设 \(A\subseteq E\)。证明
 \[
  V_A=\{f\in\mathcal B(E):f|_A=0\}
 \]
 在上确界范数下对一致极限封闭，并由正文完备性定理证明 \(V_A\) 完备。
 \exerciseitem{uniform-product-condition}
 设 \(f_n\rightrightarrows f\)、\(g_n\rightrightarrows g\)，且这些函数均
 属于 \(\mathcal B(E)\)。证明 \(f_ng_n\rightrightarrows fg\)。指出若不控制
 有界性，哪一步估计会失效。
\end{enumerate}

\subsection*{C组　构造与反例}
\begin{enumerate}[label=\textbf{\thechapter.\arabic*.},resume]
 \exerciseitem{moving-spike-integral}
 对正文的 \(p_n\) 计算图形面积，并比较峰高、支集长度与上确界范数。
 \exerciseitem{escaping-maximum-halfline}
 对 \(q_n(x)=x/(n+x)\) 证明它在每个 \([0,M]\) 上一致收敛而在
 \([0,\infty)\) 上不一致收敛。
 \exerciseitem{pointwise-unbounded-limit}
 构造一列在 \([0,1]\) 上有界的函数，使其逐点极限无界。
 \exerciseitem{uniform-convergence-preserves-infimum}
 若 \(f_n,f\in\mathcal B(E)\) 且 \(f_n\rightrightarrows f\)，证明
 \[
  \sup_E f_n\to\sup_E f,\qquad \inf_E f_n\to\inf_E f.
 \]
 \exerciseitem{uniform-convergence-varying-domain}
 令 \(E_n=[0,n]\)，\(f_n(x)=x/n\)。说明“\(f_n\) 在 \(E_n\) 上一致
 收敛”这一说法在没有共同定义域与候选极限时为什么不完整；把各函数延拓到
 \([0,\infty)\) 后给出两种不同延拓，并比较所得收敛性质。
 \optionalexerciseitem{uniform-convergence-weighted-norm}
 在 \([0,\infty)\) 上定义
 \[
  \norm f_w=\sup_{x\ge0}e^{-x}|f(x)|.
 \]
 证明它在适当的函数空间上满足范数公理，并讨论 \(f_n(x)=x/n\) 在此范数
 下是否收敛。这个例子说明改变范数会改变“共同接近”的含义。
\end{enumerate}
